\section{Tổng quan phương pháp}
Đề tài hướng tới thực thi suy luận mạng nơ-ron tại thiết bị biên bằng một hệ thống kết hợp bộ xử lý RISC-V và các bộ tăng tốc chuyên dụng. Đối tượng hiện thực là MyTinyMobileNetSoC: lõi PicoRV32 điều khiển mạng TinyMobileNet sử dụng kích hoạt và trọng số INT8. Thay vì chuyển toàn bộ đồ thị suy luận thành một khối phần cứng cố định, hệ thống phân chia nhiệm vụ theo nguyên tắc đồng thiết kế phần cứng--phần mềm (HW/SW co-design). Phần cứng thực hiện các phép nhân--cộng và tái lượng tử hóa; firmware chịu trách nhiệm thứ tự lớp, địa chỉ tensor, xử lý biên và hoàn thiện đầu ra.

Cách phân chia này cho phép giữ giao diện lập trình đơn giản qua các thanh ghi ánh xạ bộ nhớ (memory-mapped I/O, MMIO). Tuy nhiên, CPU phải chuyển toán hạng từ RAM vào từng bộ tăng tốc và đọc kết quả về. Vì vậy, thời gian suy luận phụ thuộc cả chi phí tính toán lẫn chi phí truyền dữ liệu, điều khiển vòng lặp và hậu xử lý. Đây là cơ sở để đánh giá tối ưu hóa ở cấp hệ thống trong Chương~\ref{chap:toiuu}.

Phạm vi mô tả trong chương là mã nguồn hiện tại tại \texttt{RISCVAIoT/MyTinyMobileNetSoC/}, với đường suy luận sau Level~6. Kết quả mô phỏng, kết quả hiện thực vật lý và hoạt động trên bo mạch được xem là ba loại bằng chứng riêng. Thiết kế không có DMA cho các bộ tăng tốc, không có bộ tăng tốc Softmax hoặc MaxPool, và không sử dụng kiến trúc mảng systolic làm mô hình mô tả phần cứng hiện tại. Các sơ đồ được để ở dạng vị trí chờ, nhằm tránh dùng lại hình kiến trúc không còn phù hợp.
